\section{The Treatise on Law, III: Human Law, Its Force, and Its Change (ST I-II qq.~95--97)}

\subsection{Question 95: why human law exists and where it comes from}

\textbf{The necessity of enacted law} (a.~1). Man has a natural aptitude for virtue, but its perfection requires training (\latin{disciplina}). For the well-disposed young, paternal admonition suffices. But some are found \latin{protervi et ad vitia proni}, unmoved by words, who must be restrained by force and fear---that ceasing from evil they might leave others in peace, and by habituation come at last to do willingly what they first did from dread. \latin{Huiusmodi autem disciplina cogens metu poenae, est disciplina legum. Unde necessarium fuit ad pacem hominum et virtutem, ut leges ponerentur}---this discipline that compels through fear of punishment is the discipline of laws; whence for peace and virtue it was necessary that laws be enacted; for man armed with reason can serve his lusts as no beast can, and severed from law and justice he is the worst of animals.\endnote{I-II q.~95 a.~1 co., quoting \emph{Politics} I. Note the modest, remedial account of coercive law: enacted law is necessary because of what some men are, not because obligation as such originates in enactment.}

\textbf{Derivation: the treatise's most important juridical distinction} (a.~2). Is every human law derived from the natural law? The answer opens with Augustine---\latin{non videtur esse lex, quae iusta non fuerit}, ``that which is not just seems to be no law at all''---and states the principle: in human affairs a thing is just from being right according to the rule of reason; reason's first rule is the natural law; therefore \latin{omnis lex humanitus posita intantum habet de ratione legis, inquantum a lege naturae derivatur. Si vero in aliquo, a lege naturali discordet, iam non erit lex sed legis corruptio}---every humanly enacted law has just so much of the nature of law as it is derived from the law of nature; if in anything it is discordant from the natural law, it is no longer law but a corruption of law.\endnote{I-II q.~95 a.~2 co. Section 12 returns to what this sentence does and does not license; note here only that its subject is the enactment's \emph{claim on obligation}, not a general permission of disobedience, which q.~96 a.~4 regulates separately and more cautiously.}

Then the distinction on which canonical science still runs. Something can be derived from natural law in two ways: \latin{uno modo, sicut conclusiones ex principiis; alio modo, sicut determinationes quaedam aliquorum communium}---as conclusions from principles, or as determinations of certain generalities. The first resembles demonstration: from ``harm no one'' one derives ``do not kill'' as a conclusion. The second resembles art: the builder determines the general form ``house'' to this shape; so the natural law has it that the evildoer be punished, but \emph{that he be punished with this penalty} is a determination. Both are found in human law; but conclusions have force partly from the natural law itself, while determinations \latin{ex sola lege humana vigorem habent}---have their force from human law alone.\endnote{I-II q.~95 a.~2 co.\ and ad 3: the common principles of natural law cannot be applied to all in the same way \latin{propter multam varietatem rerum humanarum}---``and hence arises the diversity of positive laws among various people,'' the treatise's calm answer to the argument from legal diversity against natural law.}

Everything a canonist means by distinguishing what the Church \emph{could} change from what she \emph{could not} is an application of this pair. A law forbidding homicide restates natural law and no legislator can repeal the obligation, only the statute. A law setting seven years as the age of canonical majority-for-obligation (c.~11) is a determination: reason requires \emph{some} threshold of discretion; nothing in nature fixes \emph{seven}; the Church fixed it and can move it.

\textbf{Isidore vindicated} (a.~4). Asked whether Isidore's division of human law is appropriate, Aquinas answers by finding its principle: a thing is divided essentially by what belongs to its notion. It belongs to the notion of human law to be derived from natural law---\latin{et secundum hoc dividitur ius positivum in ius gentium et ius civile, secundum duos modos quibus aliquid derivatur a lege naturae}: to the \latin{ius gentium} belong the conclusions---just buyings and sellings and the like, without which men cannot live together, since man is by nature a social animal; to the civil law, the particular determinations each city decides for itself.\endnote{I-II q.~95 a.~4 co. Aquinas here reads the Roman tripartition through his own conclusion/determination distinction---an interpretive act, not a report of Roman doctrine, and one that quietly reassigns the \latin{ius gentium} from ``what nations happen to share'' to ``what reason concludes.'' The article goes on to accommodate Isidore's remaining categories (military law, the forms of enactment by regime) under law's other essential notes.}

\subsection{Question 96: the force of human law}

Question 96 is the treatise's jurisprudence of limits, and its articles form a coherent doctrine of what enacted law may demand.

\textbf{Not every vice} (a.~2). Law is a rule and measure, and a measure must be homogeneous with the measured: laws must be imposed on men \latin{secundum eorum conditionem}, according to their condition---Isidore's \latin{possibilis, secundum naturam, secundum consuetudinem patriae} quoted expressly. Human law is enacted for a multitude the greater part of which is not perfect in virtue. \latin{Et ideo lege humana non prohibentur omnia vitia, a quibus virtuosi abstinent; sed solum graviora, a quibus possibile est maiorem partem multitudinis abstinere; et praecipue quae sunt in nocumentum aliorum}---human law does not forbid all vices from which the virtuous abstain, but only the graver ones from which the greater part can abstain, and chiefly those that harm others, without whose prohibition human society could not be preserved: homicide, theft, and the like.\endnote{I-II q.~96 a.~2 co.; the \latin{ad}~2 adds that law leads to virtue ``not suddenly, but gradually,'' lest precepts beyond the imperfect breed contempt and worse evils. This is not moral indifferentism but legislative prudence stated as doctrine: the scope of just prohibition is narrower than the scope of vice.}

\textbf{Not every virtue, but no virtue excluded} (a.~3). There is no virtue whose acts law cannot command---\latin{nulla virtus est de cuius actibus lex praecipere non possit}---since every virtue's objects can be referred to the common good; but human law commands only those acts \latin{ordinabiles ad bonum commune}, immediately or through the good discipline that forms citizens to preserve justice and peace.\endnote{I-II q.~96 a.~3 co.}

\textbf{Conscience, and the two kinds of unjust law} (a.~4). Do human laws bind in conscience? The article every polemicist quotes in half deserves quoting in whole. Just laws \latin{habent vim obligandi in foro conscientiae a lege aeterna, a qua derivantur}---have binding force in the court of conscience \emph{from the eternal law from which they are derived}; and laws are just from their end (the common good), their author (not exceeding his power), and their form (proportionate distribution of burdens). Unjust laws are of two kinds. Those contrary to \emph{human} good---by end (the ruler's own cupidity or vainglory), by author (beyond the power committed), or by form (unequal burdens even for the common good)---\latin{magis sunt violentiae quam leges}, are acts of violence rather than laws, since with Augustine \latin{lex esse non videtur, quae iusta non fuerit}; such laws \emph{do not bind in conscience}---\latin{nisi forte propter vitandum scandalum vel turbationem}, except perhaps to avoid scandal or disturbance, for which cause a man should yield even his right, on Matthew 5's mile and cloak. But laws unjust by contrariety to the \emph{divine} good---the laws of tyrants inducing to idolatry or to anything else against divine law---\latin{nullo modo licet observare}: may in no way be observed, \latin{quia sicut dicitur Act.\ V, obedire oportet Deo magis quam hominibus}---one must obey God rather than men.\endnote{I-II q.~96 a.~4 co. The asymmetry is the doctrine: laws unjust toward human good \emph{may} still be obeyed and sometimes should be (to avoid scandal or upheaval); laws commanding what divine law forbids \emph{may not} be obeyed at all. Between ``does not bind of itself'' and ``must be disobeyed'' lies the whole prudence of the tradition, and collapsing the two is the perennial abuse of this article.}

\textbf{Who is subject} (a.~5). Whoever is subject to a power is subject to its law; exemption arises in two ways---by being simply outside the authority (the subjects of one kingdom are not bound by another's laws), or \latin{secundum quod regitur superiori lege}, insofar as one is ruled by a higher law, as the proconsul's subject is not bound by his mandate where the emperor's governs.\endnote{I-II q.~96 a.~5 co. The higher-law exemption is the formal structure canonists still use: exemption is not lawlessness but subjection elsewhere.}

\textbf{The letter, the intention, and the birth of dispensation} (a.~6). May one subject to a law act beside its letter? Every law is ordered to the common weal and has the force of law to that extent. The legislator cannot foresee every case; he frames the law for what happens \latin{ut in pluribus}. If a case emerges in which observance would damage the common welfare---the besieged city whose law keeps the gates closed, while the enemy pursues the very citizens who defend it---the law \latin{non est observanda}: the gates are to be opened \latin{contra verba legis, ut servaretur utilitas communis, quam legislator intendit}---against the words of the law, to preserve the common benefit the legislator intended. But with a discipline attached: where the danger is not sudden, it does not belong to just anyone to decide what serves the city---\latin{hoc solum pertinet ad principes, qui propter huiusmodi casus habent auctoritatem in legibus dispensandi}, that belongs only to those in authority, who for such cases have the power of \emph{dispensing} from the laws; where the peril is too sudden for recourse, \latin{ipsa necessitas dispensationem habet annexam, quia necessitas non subditur legi}---necessity itself carries a dispensation with it, for necessity is not subject to law.\endnote{I-II q.~96 a.~6 co., quoting the Roman jurist against harsh interpretation of salutary measures; the \latin{ad}~1 insists the emergency actor ``does not judge of the law itself, but of the particular case''; the \latin{ad}~2 confines the appeal to the legislator's intention to cases of \emph{evident} harm---in doubt, follow the words or consult authority. Here, inside the treatise, is the whole later canonical apparatus in embryo: dispensation as an authority's act on its own law (CIC cc.~85--93), the emergency faculties that anticipate authority's mind (cc.~87 §2, 1079--1080), and interpretation by the mind of the legislator (c.~17)---all of it confined, note well, to \emph{human} law, since only a law framed \latin{ut in pluribus} by a legislator who cannot foresee everything has the gap this doctrine fills.}

\subsection{Question 97: change, and the force of custom}

\textbf{Why law rightly changes} (a.~1). Human law is a dictate of reason directing human acts, and may justly change on two grounds: on the part of \emph{reason}, because it is natural to human reason to advance gradually from the imperfect to the perfect---the first practical institutions, like the first philosophies, were imperfect and were corrected by successors; and on the part of \emph{men}, whose changed conditions make different things expedient---with Augustine's example of the people worthy to elect its magistrates that later, corrupted, deserves to lose the franchise.\endnote{I-II q.~97 a.~1 co., quoting \emph{De libero arbitrio} I; the \latin{ad}~1 contrasts the natural law, which \latin{immobilis perseverat}, remaining unmoved as a participation of the eternal law.}

\textbf{Why it should not change easily} (a.~2). Yet change has an intrinsic cost: \latin{ad observantiam legum plurimum valet consuetudo}---custom avails greatly for the observance of laws, so much that whatever is done against common custom, though light in itself, seems grave. When law changes, its binding power is diminished insofar as custom is taken away. Therefore human law should never be changed \latin{nisi ex aliqua parte tantum recompensetur communi saluti, quantum ex ista parte derogatur}---unless the common weal is compensated as much as it is harmed: by a very great and evident benefit, or by extreme necessity where the existing law contains manifest iniquity or its observance is very harmful.\endnote{I-II q.~97 a.~2 co., closing with the jurist's rule that in establishing novelties the benefit must be evident before departing from law long held equitable; the \latin{ad}~1 adds Aristotle: laws draw their greatest force from custom, and so are not to be quickly changed.}

\textbf{Custom's three powers} (a.~3). Can custom obtain the force of law? All law proceeds from the reason and will of the lawgiver---divine and natural law from God's reasonable will, human law from man's will regulated by reason. But reason and will are manifested by deeds as well as words: \latin{hoc enim unusquisque eligere videtur ut bonum, quod opere implet}---each man evidently chooses as good what he carries out in act. Hence by repeated acts, which make a custom, law can be changed and expounded, and something can even be established that has the force of law, \latin{inquantum per exteriores actus multiplicatos interior voluntatis motus, et rationis conceptus, efficacissime declaratur}. The conclusion, in words canon 27 of the present Code still echoes: \latin{consuetudo et habet vim legis, et legem abolet, et est legum interpretatrix}---custom has the force of law, abolishes law, and is the interpreter of laws.\endnote{I-II q.~97 a.~3 co. The \latin{ad}~3 supplies the condition the Code will make explicit: in a free multitude, the consent shown by custom counts for more than the ruler's authority; where the multitude cannot legislate for itself, custom obtains force of law only \latin{inquantum toleratur}---insofar as it is tolerated by those whose office it is to legislate, tolerance standing for approval. Compare c.~23 (section 9).}

And the ceiling, stated in the same article's first reply with a citation of Isidore: the natural and divine law proceed from the divine will, \latin{unde non potest mutari per consuetudinem procedentem a voluntate hominis \ldots\ Et inde est quod nulla consuetudo vim obtinere potest contra legem divinam vel legem naturalem}---no custom can obtain force against the divine law or the natural law; ``let custom yield to authority.''\endnote{I-II q.~97 a.~3 ad 1, quoting Isidore's \emph{Synonyma}. This sentence passes almost verbatim into c.~24 §1 of the 1983 Code (section 9): the treatise and the Code state the same rule for the same reason---human usage cannot amend what human will did not author.}

The treatise's yield can now be gathered into a map---and the map, it turns out, is not symmetrical.
